% Paper 4, §11 — Circuits: from IOI to attribution graphs
% Thesis: a circuit is a sparse, causal subgraph of the computational graph
% that explains a specific capability. Wang 2022 (IOI) is the canonical
% manual reverse-engineering; Conmy 2023 (ACDC) automates the discovery
% under causal-edge importance; Lindsey 2025 (Anthropic attribution graphs)
% scales to production via cross-layer transcoders.

\section{Circuits}
\label{sec:circuits}

A \emph{circuit} is the strongest unit of mechanistic explanation the
field has committed to: a sparse subgraph of the model's computational
graph whose nodes (attention heads, MLPs, transcoder features) and edges
(residual / attention / projection links) jointly implement an
identifiable capability. Three reference points anchor the literature.
\citet{wang2022-ioi} reverse-engineer indirect-object identification in
GPT-2 small to a 26-head subgraph in seven functional classes
\citep[\S 3]{wang2022-ioi}\footnote{KB excerpt:
\texttt{kb/excerpts/wang2022-ioi.md\#sec-3-circuit}.};
\citet{conmy2023-acdc} systematize the discovery as iterative
edge-pruning under a KL-divergence criterion\footnote{KB excerpt:
\texttt{kb/excerpts/conmy2023-acdc.md\#sec-3}.};
\citet{lindsey2025-circuit-tracing} extend the program to production
Claude 3.5 Haiku via cross-layer transcoders and Jacobian attribution.
The thesis of this section is that these are not three separate
methodologies but one: a commitment to \emph{sparse, causal, localizable
explanation}, refined and scaled across three computational regimes.

\subsection{What a ``circuit'' claims}
\label{sec:circuit-claim}

The operational definition is due to \citet{wang2022-ioi}. Let $M$ be
the model's full computational graph and $C \subseteq M$ a candidate
circuit. The function $C(x)$ is defined by \emph{knockout}: every node
in $M \setminus C$ is replaced by its mean activation on a reference
distribution $p_{\text{ABC}}$, leaving $C$'s computations
intact~\citep[\S 2]{wang2022-ioi}\footnote{KB excerpt:
\texttt{kb/excerpts/wang2022-ioi.md\#sec-2-definition}. Wang et al.\
explicitly prefer mean ablation over zero ablation: ``zero ablation
[\ldots] lead[s] to noisy results in practice'' because forced-zero
activations are off-distribution.}. A circuit claim then commits to
three quantitative criteria, also from
\citet[\S 4]{wang2022-ioi}\footnote{KB excerpt:
\texttt{kb/excerpts/wang2022-ioi.md\#sec-4-criteria}.}:
\begin{itemize}
  \item \textbf{Faithfulness:} $C$ performs the task as well as the
    whole model. Tested by metric matching on the task distribution.
  \item \textbf{Completeness:} $C$ contains every node used to perform
    the task. Tested by knocking out $M \setminus C$ and confirming
    the behavior is preserved beyond noise.
  \item \textbf{Minimality:} $C$ contains no task-irrelevant nodes.
    Tested by ablating each node in $C$ individually and verifying the
    behavior degrades.
\end{itemize}
The substantive content of a circuit claim, then, is threefold: that
the relevant features are \emph{localized} (a small fraction of the
graph), that the explanatory edges are \emph{causal} (knockout matters),
and that the subgraph is \emph{sparse} (most of $M$ is irrelevant on
the task distribution). \citet{wang2022-ioi} note that real circuits
typically pass faithfulness and minimality but fail the strongest
completeness tests~\citep[\S 4]{wang2022-ioi}; this is the crack the
\emph{backup head} phenomenon (\S\ref{sec:circuit-ioi}) widens.

\begin{intuition}
A circuit is a wiring diagram of which components matter, not a content
description of what they encode. A diagram saying ``Name Mover heads at
layer 9 attend to previous-name positions and copy them to the
unembed'' fixes the topology of the computation but leaves the question
of \emph{which direction in the residual stream} encodes ``name'' to
SAEs and probes (\S\ref{sec:sparse-autoencoders}, \S\ref{sec:probing}).
The canonical math: $C(x) = M(x)$ on a task distribution after knocking
out $M \setminus C$, with $C$ minimal under ablation.
\end{intuition}

\subsection{The IOI circuit (Wang et al.\ 2022)}
\label{sec:circuit-ioi}

The canonical end-to-end mechanistic case study. \citet{wang2022-ioi}
analyze indirect-object identification on prompts of the form ``When
\textit{Mary} and \textit{John} went to the store, \textit{John} gave a
drink to'' \(\to\) ``\textit{Mary}''. The reverse-engineered algorithm
is human-interpretable~\citep[\S 1]{wang2022-ioi}: identify all
previous names; remove duplicates; output the remainder. The circuit
that implements it consists of \textbf{26 attention heads grouped into
seven functional classes}, comprising 1.1\% of the
$(\text{head}, \text{token-position})$ pairs in GPT-2
small~\citep[Abstract]{wang2022-ioi}. The three major classes follow
the algorithm's three steps~\citep[\S 3]{wang2022-ioi}:
\begin{itemize}
  \item \textbf{Duplicate Token Heads.} Active at the second-subject
    token \texttt{S2}, attend primarily to the first-subject token
    \texttt{S1}, and write the position of the duplicate to the residual
    stream. They detect that ``John'' has been seen.
  \item \textbf{S-Inhibition Heads.} Active at the \texttt{END} token,
    attend to \texttt{S2}, and write into the \emph{query} of the Name
    Mover Heads, suppressing their attention to \texttt{S1} and
    \texttt{S2} positions.
  \item \textbf{Name Mover Heads.} Active at \texttt{END}, attend to
    previous names (now restricted to \texttt{IO} by the suppression),
    and copy the attended name to the output. The three Name Mover
    Heads have ``copy score above 95\%, compared to less than 20\% for
    an average head''~\citep[\S 3.1]{wang2022-ioi}.
\end{itemize}
Four minor classes complete the picture: \textbf{Previous Token Heads}
copy information from \texttt{S} to \texttt{S1+1} as input to the
induction route; \textbf{Induction Heads} reuse the
\citet{olsson2022-induction-heads} \texttt{[A][B]\ldots[A]\,\to\,[B]}
mechanism for duplicate-token detection rather than for in-context
learning; \textbf{Backup Name Mover Heads} replicate the Name Mover
function only when the regular Name Mover Heads are
ablated~\citep[\S 3]{wang2022-ioi}; and \textbf{Negative Name Mover
Heads} ``write in the \emph{opposite} direction of the Name Mover
Heads, thus decreasing the confidence of the
predictions''~\citep[\S 3]{wang2022-ioi}, which the authors speculate
is loss-hedging behavior.

The discovery method is \emph{path patching}~\citep[\S 3.1]{wang2022-ioi}.
For inputs $x_{\text{orig}}$ and $x_{\text{new}}$ and a candidate head
$h$, path patching ``runs a forward pass on $x_{\text{orig}}$, but for
the paths in $\mathcal{P}$ it replaces the activations for $h$ with
those from $x_{\text{new}}$''~\citep[\S 3.1]{wang2022-ioi}, where
$\mathcal{P}$ is the set of direct paths from $h$ to the unembedding
through residual connections and MLPs but \emph{not} through other
attention heads. This isolates the direct effect of $h$ on the logits
from indirect routing through other heads --- the partial-derivative
analogue of full activation patching, and the methodological refinement
that makes the seven-class decomposition legible.

The IOI circuit complicated three field expectations
simultaneously~\citep[\S 3]{wang2022-ioi}:
(i) \emph{redundancy} --- ablating Name Movers does not kill the
behavior, because Backup Name Movers latently reactivate; circuits are
not rigid pipelines.
(ii) \emph{repurposed mechanism} --- induction heads, originally
identified for in-context learning, are reused here for duplicate
detection; mechanism does not equal function.
(iii) \emph{negative components} --- a head that consistently writes
\emph{against} the correct answer is a structural feature of the
trained model, not a bug. Each finding shows up again at scale in the
attribution-graph era (\S\ref{sec:circuit-attribution}).

\subsection{Algorithmic circuit discovery (Conmy et al.\ 2023)}
\label{sec:circuit-acdc}

Hand-discovering the IOI circuit took multiple researcher-months.
\citet{conmy2023-acdc} systematize this work as an algorithm
(\textbf{ACDC}, Automatic Circuit DisCovery) over an explicitly
constructed computational DAG. The graph $G$ represents ``interactions
between attention heads and MLPs''~\citep[\S 2.2]{conmy2023-acdc} ---
or, at finer granularity, between separate query/key/value activations,
individual neurons, or per-token-position
nodes\footnote{KB excerpt:
\texttt{kb/excerpts/conmy2023-acdc.md\#sec-2-2}.}. The
\emph{interchange-intervention} corruption convention --- preferred by
\citet{conmy2023-acdc} over zero or mean ablation because both ``take
the model too far away from actually possible activation
distributions''~\citep[\S 2.3]{conmy2023-acdc} --- fixes the corrupted
input $x'_i$ as a same-distribution counterfactual drawn from a
related prompt template.

The metric for edge importance is \emph{KL-divergence} between the full
graph's output distribution $G(x_i)$ and the candidate subgraph's
output distribution $H(x_i, x'_i)$, where $H(x_i, x'_i)$ is computed by
``overwrit[ing] all edges in $G$ that are not present in $H$ to their
activation on $x'_i$''~\citep[\S 3]{conmy2023-acdc}. ACDC then iterates
the graph in reverse-topological order (output-to-input), at each node
greedily attempting to remove every incoming edge whose removal does
not increase the KL divergence by more than threshold $\tau$. Verbatim
from~\citet[Algorithm 1]{conmy2023-acdc}\footnote{KB excerpt:
\texttt{kb/excerpts/conmy2023-acdc.md\#sec-3}.}:

\begin{lstlisting}[caption={ACDC, the Automatic Circuit DisCovery algorithm.},label={lst:acdc}]
Algorithm 1: ACDC
Data:   Computational graph G; dataset (x_i)_{i=1..n};
        corrupted datapoints (x'_i)_{i=1..n}; threshold tau > 0.
Result: Subgraph H subset of G.

1  H <- G                              // initialize H to the full graph
2  H <- H.reverse_topological_sort()   // sort H so output node first
3  for v in H do
4    for w parent of v do
5      H_new <- H \ {w -> v}           // tentatively remove candidate edge
6      if D_KL(G || H_new) - D_KL(G || H) < tau then
7        H <- H_new                    // edge unimportant, prune permanently
8      end
9    end
10 end
11 return H
\end{lstlisting}

The algorithm's output is a sparse subgraph $H \subseteq G$ that
preserves the KL-divergence to the full model up to $\tau$ summed over
prunings. On GPT-2 small Greater-Than, ACDC ``rediscovered 5/5 of the
component types in [the manually-found] circuit'' and ``selected 68 of
the 32{,}000 edges in GPT-2 Small, all of which were manually found by
previous work''~\citep[Abstract]{conmy2023-acdc}\footnote{KB excerpt:
\texttt{kb/excerpts/conmy2023-acdc.md\#abstract}.}. ACDC is evaluated
against two baselines: \emph{Subnetwork Probing} (a learned mask over
components, with the linear probe removed and KL substituted for the
training metric) and \emph{Head Importance Score for Pruning} (a
top-$k$ ranking by head importance score). Both are generalized in
\citet{conmy2023-acdc} to use corrupted-activation patching rather
than zero patching, for fair comparison. ROC-curve analysis against
canonical circuits answers the two evaluation questions
\citep[\S 4]{conmy2023-acdc}: does the method find the right edges
(true-positive rate) and does it avoid spurious ones (false-positive
rate)?

ACDC is the algorithmic spine of automated circuit discovery, but it
inherits the cost of activation patching: each candidate edge requires
a forward pass with one substituted activation, so the per-task cost
scales with the size of the computational graph. On GPT-2-small that
is feasible; on production models it is not.

\subsection{Attribution graphs (Anthropic 2025)}
\label{sec:circuit-attribution}

\citet{lindsey2025-circuit-tracing} synthesize the program for
production-scale models. The construction has three pieces. First,
\emph{cross-layer transcoders} replace each MLP sublayer with a sparse
factorization whose decoder writes into the residual stream of every
\emph{subsequent} layer, not just the immediate next layer; this makes
skip-layer feature flow explicit in the
graph~\citep{dunefsky2024-transcoders}. Second, a single forward pass
on the replaced model records the activations of all transcoder
features $\{f_i^{(\ell, t)}\}$. Third, an \emph{attribution graph} is
constructed by computing the linear contribution of every active
feature to a target node (typically the unembedding direction for
the target token), via the Jacobian of the cross-layer transcoder
system.

Concretely, the attribution edge from feature $i$ at layer $\ell_1$,
position $t_1$ to feature $j$ at layer $\ell_2$, position $t_2$ is
\begin{equation}
  a\!\left(f_i^{(\ell_1, t_1)} \to f_j^{(\ell_2, t_2)}\right)
  \;=\;
  \frac{\partial f_j^{(\ell_2, t_2)}}{\partial f_i^{(\ell_1, t_1)}}
  \;\bigg|_{x}
  \;\;\approx\;\;
  \big\langle \mathbf{e}_j^{(\ell_2)},
              \mathbf{A}_{\ell_1 \to \ell_2}\,
              \mathbf{d}_i^{(\ell_1)} \big\rangle,
  \label{eq:attribution-edge}
\end{equation}
where $\mathbf{d}_i^{(\ell_1)}$ is the decoder column of feature $i$
in the layer-$\ell_1$ transcoder, $\mathbf{e}_j^{(\ell_2)}$ is the
encoder row of feature $j$ in the layer-$\ell_2$ transcoder, and
$\mathbf{A}_{\ell_1 \to \ell_2}$ accumulates the linearized residual /
attention transport between the two layers (identity through the
residual stream; attention-weight-modulated through any intervening
attention layers). The bilinear factorization is what makes
Eq.~\eqref{eq:attribution-edge} input-invariant in its endpoints and
input-dependent only through the attention weights and the active set
$\{f_i^{(\ell, t)} > 0\}$. Edges are pruned below threshold; the
resulting graph is iterated backward from the target node to the
embedding layer; the surviving subgraph is the attribution
graph~\citep{lindsey2025-circuit-tracing}.

\begin{forumsignal}
\textbf{Single-source methodology (2026-Q2).} The cross-layer-transcoder
+ Jacobian-attribution pipeline above rests on
\citet{lindsey2025-circuit-tracing} as its sole canonical reference; no
external lab has published an independent reproduction at frontier
scale at this writing. Treat the productionised-methodology framing as
Anthropic-internal evidence pending external replication. The closing
contradiction of this section develops the attendant
transcoder-seed-stability and external-reproduction concerns.
\end{forumsignal}

The compute regime shifts from manual to automated. Per-task cost on
the replaced model is one forward pass plus one backward pass; the
amortized cost is the cross-layer transcoder training, which Anthropic
estimates at roughly comparable order to SAE training (Gemma Scope
reports SAE training at $\sim$20\% of the underlying model's
pretraining compute~\citep{gemma-scope-2024}). Reported findings on
Claude 3.5 Haiku, transferred from
\citet{lindsey2025-biology}\deepencite{lindsey2025-biology HTML-only;
verbatim quotes pending Phase 2 transcription}: forward-planning
features that activate many tokens before a poetic rhyme is produced;
multi-step factual chains in which a ``Texas'' feature in middle layers
licenses an ``Austin'' feature in later layers for the prompt ``the
capital of the state Dallas is in''; refusal-circuit features on
harmful-request prompt classes; and partial cross-lingual feature
overlap with language-routing through a partly language-agnostic middle.

\begin{intuition}
ACDC and attribution graphs answer the same question --- which edges of
the computational graph causally implement a behavior --- in two
compute regimes. ACDC pays $O(|E|)$ forward passes per task and
returns components; attribution graphs amortize transcoder training
once and pay one forward+backward per task, returning features. The
canonical math: ACDC prunes by
$D_{\mathrm{KL}}(G \| H_{\text{new}}) - D_{\mathrm{KL}}(G \| H) < \tau$
(\S\ref{sec:circuit-acdc}); attribution graphs prune by the linearized
Jacobian in Eq.~\eqref{eq:attribution-edge}.
\end{intuition}

\subsection{When circuits agree across methods}
\label{sec:circuit-agreement}

The IOI circuit is the canonical convergence story.
\citet{wang2022-ioi}'s manual path-patching reverse-engineering and
\citet{conmy2023-acdc}'s ACDC recover overlapping subgraphs --- ACDC
``rediscovered 5/5 of the component types''~\citep[Abstract]{conmy2023-acdc}
on related GPT-2-small tasks, and on IOI itself recovers the seven
functional classes Wang et al.\ describe. SAE-feature analysis on the
same prompts identifies features playing the QK and OV roles of the
discovered heads (\S\ref{sec:sparse-autoencoders}). Tuned-lens
prediction trajectories on IOI prompts are consistent with the
circuit's mid-late-layer commit to the IO token
(\S\ref{sec:logit-lens-family}). Three independent methodologies
converge on one picture.

Factual recall is the second convergence. \citet{meng2022-rome}'s
causal tracing on GPT-2 XL localizes factual recall to mid-layer MLPs
at the last subject token (AIE peaks at layer 15;
$\mathrm{AIE}_{\mathrm{MLP}} = 6.6\%$ vs.\
$\mathrm{AIE}_{\mathrm{attn}} = 1.6\%$
\citep[\S 2.2]{meng2022-rome}, with rank-one weight edits at exactly
that locus achieving 99.8\% efficacy on
zsRE~\citep[\S 3.2]{meng2022-rome}). SAE features on those layers
include subject-entity features
\citep{templeton2024-scaling-monosemanticity}\deepencite{templeton2024-scaling-monosemanticity
§3 Haiku features --- HTML-only, awaiting transcription}; the tuned
lens shows the predicted-object token's probability rising across the
localized layer band; ROME's rank-one MLP edit succeeds at exactly
that locus. Four methods, one converging picture.

The induction-head story (\citet{olsson2022-induction-heads}) is a
third. The same two-attention-head pattern is identified by ablation
(removing the heads abolishes in-context learning), by phase-change
correlation (induction-head formation aligns with the in-context
learning loss-curve discontinuity), by direct circuit
reverse-engineering, and is then \emph{reused} as a substructure of
the IOI circuit~\citep[\S 3]{wang2022-ioi}\deepencite{olsson2022-induction-heads
HTML-only at transformer-circuits.pub --- six lines of evidence
narrative pending Phase 2 transcription}. The cross-method evidence is
strong enough that ``induction head'' is now a load-bearing primitive
across the rest of the literature.

\subsection{When circuits disagree}
\label{sec:circuit-disagreement}

Convergence is not universal. The truth-direction circuit
(\S\ref{sec:truth-directions}) is the cleanest divergence case. The
mass-mean direction $\mathbf{d}_{\mathrm{MM}} = \boldsymbol{\mu}_+ -
\boldsymbol{\mu}_-$ found by \citet{marks-tegmark-2023-truth} is
dataset-dependent in early-to-mid layers --- antipodal between
\texttt{cities} and \texttt{neg\_cities}, then orthogonal, only aligned
across formats in late layers~\citep[\S 4]{marks-tegmark-2023-truth}%
\footnote{KB excerpt:
\texttt{kb/excerpts/marks-tegmark-2023-truth.md\#sec-4-misalignment}.}.
There is no single ``truth circuit'' in this picture: there is a
\emph{family} of dataset-conditional truth representations, partially
aligned across formats and partially not, with the alignment depth
itself varying across model layers.

Sycophancy is a parallel non-convergence. Probing-derived steering
directions for sycophancy are reported across multiple labs, but no
two reports agree on a single ``sycophancy circuit'': the proposed
mechanisms span steering-vector additions to refusal-feature
ablations to fine-tuning interventions, and the published evidence
does not yet localize a stable subgraph. The 2025 finding that
sycophancy decomposes into agreement, praise, and genuine-agreement
components individually suggests \emph{behaviorally singular,
representationally plural} --- if confirmed, the question ``what is
the sycophancy circuit?'' is malformed: there are several.

\begin{contradiction}
\textbf{Circuit-tracing scalability past Anthropic-2025.} The
attribution-graph methodology of \citet{lindsey2025-circuit-tracing}
and the Haiku findings of \citet{lindsey2025-biology} have not been
independently reproduced as of 2026. The compute cost
($\sim\!20\%$ of pretraining for transcoder training) is a real
adoption barrier outside frontier labs. The substantive claims
--- forward planning, multi-step factual chains, refusal circuits ---
rest on internal Anthropic evidence and on transcoders whose
canonicality inherits the \citet{leask2025-sae-not-canonical}
critique: a different transcoder training run may factor the same MLP
into a different sparse basis, and the stability of attribution-graph
conclusions across transcoder seeds is open empirical. The strongest
form of the claim --- that automated circuit discovery scales to
Claude-class models with results that converge across replications
--- is asserted but not yet externally validated.
\end{contradiction}

\subsection{Forward link}
\label{sec:circuit-forward}

The cases above set up \S\ref{sec:cross-method}, where we make the
convergence/divergence pattern the object of analysis: when methods
should agree (because they answer the same question about the same
mechanism), divergence falsifies at least one method's assumptions;
when methods answer different questions (correlational vs.\ causal,
component-level vs.\ feature-level), apparent divergence may be
principled. The IOI circuit and factual-recall localization are the
positive cases; the truth-direction geometry and SAE canonicality are
the negative cases; and the open question --- whether attribution-graph
findings on Haiku will reproduce outside Anthropic --- is the load-bearing
test for whether circuit-tracing extends from a method to a paradigm.
